\documentclass[8pt,a4paper]{article}
\usepackage[utf8]{inputenc}
\usepackage[italian]{babel}
\usepackage{graphicx}
\usepackage{amsmath}
\usepackage{float}
\usepackage{listings}
\usepackage{xcolor}
\usepackage{hyperref}
\usepackage{booktabs}
\usepackage{geometry}

\geometry{margin=0.65in}

\lstset{
    language=Scala,
    basicstyle=\ttfamily\small,
    keywordstyle=\color{blue},
    commentstyle=\color{gray},
    breaklines=true,
    frame=single,
    numbers=left
}

\title{Earthquake Co-occurrence Analysis: Strong Scalability Study}
\date{\today}

\begin{document}
\maketitle

\section{Introduzione}
Questo progetto implementa un'applicazione in Scala per l'analisi di co-occorrenza di eventi sismici utilizzando Spark su Google Cloud Platform (Dataproc). L'obiettivo è identificare coppie di locazioni geografiche che sperimentano terremoti nello stesso giorno e valutare la scalabilità dell'applicazione su cluster di diverse dimensioni. Sono state testate due configurazioni di macchine virtuali (n2-standard-4 e n1-standard-2) con benchmarking automatizzato tramite script PowerShell.

\section{Setup Infrastruttura Google Cloud}

Prima di eseguire il benchmarking, è stato necessario configurare l'infrastruttura su Google Cloud Platform tramite lo script PowerShell \texttt{setup\_gcloud.ps1}, che automatizza i seguenti compiti:

\begin{enumerate}
    \item \textbf{Configurazione Progetto GCP}: imposta il progetto \texttt{scalableproject-482714}
    \item \textbf{Creazione Bucket Cloud Storage}: crea \texttt{terremoti-bucket-giorgiapirelli} con cartelle dedicate (\texttt{jars/}, \texttt{output/})
    \item \textbf{Build Automatico con SBT}: compila con \texttt{sbt clean} e \texttt{sbt package}
    \item \textbf{Upload JAR in Cloud Storage}: carica il JAR compilato
    \item \textbf{Verifica Dataset}: Controlla la presenza del dataset
\end{enumerate}

Lo script deve essere eseguito una sola volta prima del benchmarking.

\section{Workflow del Sistema}

Il workflow viene sviluppato:

\begin{enumerate}
    \item \textbf{Creazione Cluster}: script PowerShell crea cluster Dataproc con nodo master + N worker
    \item \textbf{Esecuzione Job}: lancia l'applicazione via \texttt{gcloud dataproc jobs submit spark}
    \item \textbf{Misurazione Performance}: registra tempi di esecuzione in CSV
    \item \textbf{Pulizia}: elimina il cluster automaticamente dopo il job
\end{enumerate}

\section{Approccio Implementativo}

L'applicazione implementa quattro fasi di elaborazione:

\subsection{Fase 1: Parsing e Trasformazione}
Estrazione di (data, latitudine, longitudine) da CSV, filtro header, arrotondamento coordinate alla prima cifra decimale.

\subsection{Fase 2: Partizionamento}
Utilizzo di HashPartitioner per distribuzione uniforme per data:
\begin{lstlisting}
val partitioned = data.partitionBy(new HashPartitioner(numPartitions))
val grouped = partitioned.groupByKey()
\end{lstlisting}

Garantisce: distribuzione uniforme.

\subsection{Fase 3: Generazione Coppie e Ricerca del Massimo}

Per ogni data, vengono generate tutte le possibili coppie di locazioni co-occorrenti. Le fasi specifiche sono:

\begin{itemize}
    \item \textbf{Generazione}: tramite \texttt{flatMap}, vengono create tutte le combinazioni non ordinate di locazioni per ogni data
    \item \textbf{Deduplicazione}: le locazioni vengono trasformate in Set per eliminare automaticamente i duplicati
    \item \textbf{Raggruppamento}: con \texttt{groupByKey()}, tutte le date associate a ogni coppia vengono raccolte insieme
    \item \textbf{Ricerca del Massimo}: mediante l'operazione \texttt{reduce()}, viene identificata la coppia che co-occorrerà nel maggior numero di date distinte
\end{itemize}

\subsection{Fase 4: Output}
Salvataggio coppia vincente e lista di date di co-occorrenza.

\section{Configurazioni di Test}

Durante i test iniziali con n2-standard-4, si è riscontrato un errore di quota GCP nella creazione del cluster con 4 nodi worker, causando il fallimento della creazione del cluster con tale configurazione.

\begin{figure}[H]
\centering
\includegraphics[width=0.9\textwidth]{error.png}
\caption{errore di quota GCP, 20 CPU richiesti vs 16 CPU disponibili nel progetto}
\label{fig:error_quota}
\end{figure}
\noindent
A causa di questo vincolo di quota GCP, è stata adottata una strategia di test alternativa e più completa:

\begin{enumerate}
    \item \textbf{Test con n2-standard-4}: limitati a 2-3 worker con dataset completo per rispettare il vincolo di quota
    \item \textbf{Test con n1-standard-2}: utilizzata come configurazione alternativa con 2 vCPU, 7.5 GB RAM, 100 GB disk per nodo. Questa scelta ha permesso test più estesi:
    \begin{itemize}
        \item \textbf{Dataset piccolo}: 2-4 worker con 4/8/12 partizioni
        \item \textbf{Dataset completo}: 2-4 worker con 16/24 partizioni
    \end{itemize}
\end{enumerate}
\noindent
Sebbene il progetto richiedesse l'utilizzo del solo dataset completo, si è scelto di includere anche i risultati per il dataset piccolo con n1-standard-2 per fornire un'analisi comparativa più completa. 

\subsection{Motivazione dei test con più partizioni}

I vari test e prove sono stati eseguiti aumentando progressivamente il numero di partizioni perché \textbf{si riduceva il tempo di esecuzione}. Invece, se si manteneva troppo basso il numero di partizioni, la riduzione dei tempi \textbf{non era significativa}. 
\noindent
Per questo motivo, i test sono stati eseguiti partendo da \textbf{poche partizioni (4)} fino a \textbf{molte partizioni (24)} per esplorare sistematicamente l'intero range.

Esempi dal range completo testato:

\begin{itemize}
    \item \textbf{n2-standard-4 con 2 worker} (range 8-24):
    \begin{itemize}
        \item 8 partizioni: 923.05s
        \item 16 partizioni: 641.63s (-30.4\%, miglioramento)
        \item 24 partizioni: 397.42s (-38.1\%, continua a migliorare)
    \end{itemize}
    
    \item \textbf{n1-standard-2 dataset completo con 2 worker} (range 16-24):
    \begin{itemize}
        \item 16 partizioni: 2412.54
        \item 24 partizioni: 888.59s (-63.1\%, miglioramento significativo)
    \end{itemize}
    
    \item \textbf{n1-standard-2 dataset piccolo con 2 worker} (range 4-12):
    \begin{itemize}
        \item 4 partizioni: 1904.07s 
        \item 8 partizioni: 1172.29s (-38.4\%, miglioramento)
        \item 12 partizioni: 523.30s (-55.4\%, continua a migliorare)
    \end{itemize}
\end{itemize}
\noindent
Questo approccio sistematico ha permesso di scoprire che aumentare progressivamente le partizioni produceva consistenti riduzioni dei tempi di esecuzione, giustificando l'esecuzione di test con il range completo (da 4 a 24 partizioni).


\section{Metodologia di calcolo: strong scalability su GCP Dataproc}

\subsection{Formula standard vs Formula GCP Dataproc}

La metrica di efficienza di scalabilità è definita come:

\begin{equation}
E(n) = \frac{S(n)}{n} = \frac{T(1)}{n \times T(n)}
\end{equation}
\noindent
\textbf{Teoricamente}, $E(n) \leq 100\%$ rappresenta la scalabilità lineare perfetta. Tuttavia, il valore $T(1)$ non è disponibile sperimentalmente.

\subsection{Soluzione: Formula Alternativa per GCP Dataproc}

Quindi si utilizza $T(2)$ (tempo con 2 worker) come baseline, poiché è il minimo misurabile. La formula diventa:

\begin{equation}
E(n) = \frac{S(n)}{n} \times 2 = \frac{T(2)}{n \times T(n)} \times 2 = \frac{2}{n} \times S(n)
\end{equation}

dove $S(n) = T(2) / T(n)$.
\textbf{Esempio Numerico}: con 3 worker: $E(3) = (2/3) \times 1.62 \times 100 = 108.0\%$

\section{Analisi di Scalabilità}

\subsection{Risultati n1-standard-2 (dataset piccolo)}

\begin{table}[H]
\centering
\caption{Benchmark n1-standard-2 - dataset piccolo}
\label{tab:n1_results}
\footnotesize
\begin{tabular}{|c|c|c|c|c|}
\hline
\textbf{W} & \textbf{P} & \textbf{Durata (s)} & \textbf{Speedup} & \textbf{Eff} \\
\hline
2 & 4 & 1904.07 & 1.00 & 100.0\% \\
2 & 8 & 1172.29 & 1.00 & 100.0\% \\
2 & 12 & 523.30 & 1.00 & 100.0\% \\
\hline
3 & 4 & 1429.39 & 1.33 & 88.7\% \\
3 & 8 & 753.40 & 1.56 & 104.0\% \\
3 & 12 & 431.61 & 1.21 & 80.7\% \\
\hline
4 & 4 & 995.01 & 1.91 & 95.5\% \\
4 & 8 & 708.06 & 1.65 & 82.5\% \\
4 & 12 & 349.71 & 1.50 & 75.0\% \\
\hline
\end{tabular}
\end{table}

\subsection{Risultati n1-standard-2 (dataset completo)}

I test con n1-standard-2 sono stati estesi al dataset completo.

\begin{table}[H]
\centering
\caption{Benchmark n1-standard-2 - dataset completo}
\label{tab:n1_full_results}
\footnotesize
\begin{tabular}{|c|c|c|c|c|c|}
\hline
\textbf{W} & \textbf{P} & \textbf{Durata (s)} & \textbf{Speedup} & \textbf{Eff} \\
\hline
2 & 16 & 2412.54 & 1.00 & 100.0\% \\
2 & 24 & 888.59 & 1.00 & 100.0\% \\
\hline
3 & 16 & 1489.71 & 1.62 & 108.0\% \\
3 & 24 & 640.86 & 1.39 & 92.7\% \\
\hline
4 & 16 & 1109.71 & 2.17 & 108.5\% \\
4 & 24 & 492.26 & 1.81 & 90.5\% \\
\hline
\end{tabular}
\end{table}


\subsubsection{Analisi dettagliata per numero di partizioni - n1-standard-2 dataset completo}

\begin{table}[H]
\centering
\caption{Scalabilità con 16 partizioni (n1-standard-2 - dataset completo)}
\begin{tabular}{|c|c|c|c|}
\hline
\textbf{Workers} & \textbf{Durata (s)} & \textbf{Speedup} & \textbf{Efficienza} \\
\hline
2 & 2412.54 & 1.00 & 100.0\% \\
3 & 1489.71 & 1.62 & 108.0\% \\
4 & 1109.71 & 2.17 & 108.5\% \\
\hline
\end{tabular}
\end{table}

\begin{table}[H]
\centering
\caption{Scalabilità con 24 partizioni (n1-standard-2 - dataset completo)}
\begin{tabular}{|c|c|c|c|}
\hline
\textbf{Workers} & \textbf{Durata (s)} & \textbf{Speedup} & \textbf{Efficienza} \\
\hline
2 & 888.59 & 1.00 & 100.0\% \\
3 & 640.86 & 1.39 & 92.7\% \\
4 & 492.26 & 1.81 & 90.5\% \\
\hline
\end{tabular}
\end{table}

\subsection{Risultati n2-standard-4 (dataset completo)}

Questo set di test con n2-standard-4 e dataset completo mostra i risultati con configurazioni di partizioni da 8 a 24.

\begin{table}[H]
\centering
\caption{Benchmark complessivo n2-standard-4 - Tutte le partizioni}
\label{tab:n2_complete_all}
\footnotesize
\begin{tabular}{|c|c|c|c|c|}
\hline
\textbf{Workers} & \textbf{Partizioni} & \textbf{Durata (s)} & \textbf{Speedup} & \textbf{Efficienza} \\
\hline
2 & 8 & 911.99 & 1.00 & 100.0\% \\
2 & 16 & 537.13 & 1.00 & 100.0\% \\
2 & 24 & 389.19 & 1.00 & 100.0\% \\
\hline
3 & 8 & 772.00 & 1.18 & 78.7\% \\
3 & 16 & 450.10 & 1.19 & 79.3\% \\
3 & 24 & 284.99 & 1.37 & 91.3\% \\
\hline
\end{tabular}
\end{table}

\textbf{Osservazioni n2-standard-4 completo:} configurazione ottimale con 24 partizioni + 3 worker (speedup 1.37x, efficienza 91.3\%)

\subsubsection{Analisi dettagliata per partizioni - n2-standard-4 dataset completo}

\begin{table}[H]
\centering
\caption{Scalabilità con 8 partizioni (n2-standard-4 - Dataset Completo)}
\begin{tabular}{|c|c|c|c|}
\hline
\textbf{Workers} & \textbf{Durata (s)} & \textbf{Speedup} & \textbf{Efficienza} \\
\hline
2 & 911.99 & 1.00 & 100.0\% \\
3 & 772.00 & 1.18 & 78.7\% \\
\hline
\end{tabular}
\end{table}

\begin{table}[H]
\centering
\caption{Scalabilità con 16 partizioni (n2-standard-4 - Dataset Completo)}
\begin{tabular}{|c|c|c|c|}
\hline
\textbf{Workers} & \textbf{Durata (s)} & \textbf{Speedup} & \textbf{Efficienza} \\
\hline
2 & 537.13 & 1.00 & 100.0\% \\
3 & 450.10 & 1.19 & 79.3\% \\
\hline
\end{tabular}
\end{table}

\begin{table}[H]
\centering
\caption{Scalabilità con 24 partizioni (n2-standard-4 - Dataset Completo)}
\begin{tabular}{|c|c|c|c|}
\hline
\textbf{Workers} & \textbf{Durata (s)} & \textbf{Speedup} & \textbf{Efficienza} \\
\hline
2 & 389.19 & 1.00 & 100.0\% \\
3 & 284.99 & 1.37 & 91.3\% \\
\hline
\end{tabular}
\end{table}


\subsubsection{Analisi Dettagliata per Numero di Partizioni - n2-standard-4}

\begin{table}[H]
\centering
\caption{Scalabilità con 8 partizioni (n2-standard-4)}
\begin{tabular}{|c|c|c|c|}
\hline
\textbf{Workers} & \textbf{Durata (s)} & \textbf{Speedup} & \textbf{Efficienza} \\
\hline
2 & 923.05 & 1.00 & 100.0\% \\
3 & 883.41 & 1.04 & 69.3\% \\
\hline
\end{tabular}
\end{table}

\begin{table}[H]
\centering
\caption{Scalabilità con 16 partizioni (n2-standard-4)}
\begin{tabular}{|c|c|c|c|}
\hline
\textbf{Workers} & \textbf{Durata (s)} & \textbf{Speedup} & \textbf{Efficienza} \\
\hline
2 & 641.63 & 1.00 & 100.0\% \\
3 & 526.52 & 1.22 & 81.3\% \\
\hline
\end{tabular}
\end{table}

\begin{table}[H]
\centering
\caption{Scalabilità con 24 partizioni (n2-standard-4)}
\begin{tabular}{|c|c|c|c|}
\hline
\textbf{Workers} & \textbf{Durata (s)} & \textbf{Speedup} & \textbf{Efficienza} \\
\hline
2 & 397.42 & 1.00 & 100.0\% \\
3 & 288.43 & 1.38 & 92.0\% \\
\hline
\end{tabular}
\end{table}


\subsection{Seconda prova con n2-standard-4}

Questo secondo set di test con n2-standard-4 e dataset completo mostra i risultati con configurazioni di partizioni da 4 a 24.
\begin{table}[H]
\centering
\caption{Benchmark completo n2-standard-4 - Tutte le partizioni}
\label{tab:n2_complete_all}
\footnotesize
\begin{tabular}{|c|c|c|c|}
\hline
\textbf{Workers} & \textbf{Durata (s)} & \textbf{Speedup} & \textbf{Efficienza} \\
\hline
2 & 1034.07 & 1.00 & 100.0\% \\
2 & 920.30 & 1.00 & 100.0\% \\
2 & 583.71 & 1.00 & 100.0\% \\
2 & 400.04 & 1.00 & 100.0\% \\
\hline
3 & 862.90 & 1.20 & 80.0\% \\
3 & 773.76 & 1.19 & 79.3\% \\
3 & 380.04 & 1.54 & 102.7\% \\
3 & 275.39 & 1.45 & 96.7\% \\
\hline
\end{tabular}
\end{table}
\noindent
\textbf{Osservazioni n2-standard-4 completo:} configurazione ottimale con 16 partizioni + 3 worker (speedup 1.54x, efficienza 102.7\%)

\subsubsection{Analisi Dettagliata per Partizioni - n2-standard-4 completo}

\begin{table}[H]
\centering
\caption{Scalabilità con 4 partizioni (n2-standard-4)}
\begin{tabular}{|c|c|c|c|}
\hline
\textbf{Workers} & \textbf{Durata (s)} & \textbf{Speedup} & \textbf{Efficienza} \\
\hline
2 & 1034.07 & 1.00 & 100.0\% \\
3 & 862.90 & 1.20 & 80.0\% \\
\hline
\end{tabular}
\end{table}

\begin{table}[H]
\centering
\caption{Scalabilità con 8 partizioni (n2-standard-4)}
\begin{tabular}{|c|c|c|c|}
\hline
\textbf{Workers} & \textbf{Durata (s)} & \textbf{Speedup} & \textbf{Efficienza} \\
\hline
2 & 920.30 & 1.00 & 100.0\% \\
3 & 773.76 & 1.19 & 79.3\% \\
\hline
\end{tabular}
\end{table}

\begin{table}[H]
\centering
\caption{Scalabilità con 16 partizioni (n2-standard-4)}
\begin{tabular}{|c|c|c|c|}
\hline
\textbf{Workers} & \textbf{Durata (s)} & \textbf{Speedup} & \textbf{Efficienza} \\
\hline
2 & 583.71 & 1.00 & 100.0\% \\
3 & 380.04 & 1.54 & 102.7\% \\
\hline
\end{tabular}
\end{table}

\begin{table}[H]
\centering
\caption{Scalabilità con 24 partizioni (n2-standard-4)}
\begin{tabular}{|c|c|c|c|}
\hline
\textbf{Workers} & \textbf{Durata (s)} & \textbf{Speedup} & \textbf{Efficienza} \\
\hline
2 & 400.04 & 1.00 & 100.0\% \\
3 & 275.39 & 1.45 & 96.7\% \\
\hline
\end{tabular}
\end{table}

\subsubsection{Analisi Comparativa: Tutte le Configurazioni}

La Tabella \ref{tab:comparison_all} confronta i risultati delle configurazioni ottimali tra dataset piccolo n1-standard-2, dataset completo n1-standard-2, e dataset completo n2-standard-4.

\begin{table}[H]
\centering
\caption{Confronto configurazioni ottimali}
\label{tab:comparison_all}
\begin{tabular}{|l|l|c|c|c|c|}
\hline
\textbf{Machine} & \textbf{Dataset} & \textbf{Config Opt} & \textbf{Speedup} & \textbf{Efficienza} \\
\hline
n1-standard-2 & Piccolo  & 8p + 3w & 1.56x & 104.0\% \\
n1-standard-2 & Completo & 16p + 4w & 2.17x & 108.5\% \\
n2-standard-4 & Completo & 24p + 3w & 1.38x & 92.0\% \\
\hline
\end{tabular}
\end{table}

\section{Codice Sorgente:} \url{https://github.com/giorgiapirelli/ScalableProjectGiorgiaPirelli.git}

\end{document}